%% kapa: exam=1 round=23 year=2012 period=2 form=A subject=accounting subject_ko=회계학 seq=19 number=19
%% source: Q-Net 감정평가사 기출문제 (cst00309) artlSeq=5052580 file=raw/1008/5052580/2012년 제23회 감정평가사 제1차시험 기출문제(2교시 A형).pdf
%% conversion: kordoc 4.24.51 -> tools/gen1.py (auto); review=proofread-sonnet; flags=none
\begin{question}{19}
(주)서울은 20×1년 1월 1일 ￦1,000,000에 기계장치를 취득하여 사용하기 시작하였다. 기계장치의 내용연수는 5년이고 잔존가치 없이 정액법으로 상각하며 원가모형을 적용한다. (주)서울은 20×2년부터 기계장치에 대해서 재평가모형을 최초 적용하기로 하였다. 또한 내용연수를 재검토한 결과 당초 내용연수를 5년이 아니라 6년으로 재추정하였다. 20×2년 12월 31일 기계장치의 공정가치는 ￦700,000이다. (주)서울이 20×2년에 인식할 기계장치의 감가상각비와 20×2년 12월 31일 재평가잉여금의 잔액은?
\begin{choices}
  \item 감가상각비 ￦133,333 재평가잉여금 ￦33,333
  \item 감가상각비 ￦133,333 재평가잉여금 ￦0
  \item 감가상각비 ￦166,667 재평가잉여금 ￦66,667
  \item 감가상각비 ￦160,000 재평가잉여금 ￦60,000
  \item 감가상각비 ￦160,000 재평가잉여금 ￦0
\end{choices}
\end{question}
\answer{4}
